\documentclass[10pt,a4paper]{amsart}
\usepackage[margin=19mm]{geometry}
\usepackage{amsmath,amssymb,mathtools}
\usepackage[scr=rsfs,cal=euler]{mathalfa}
\usepackage{xfrac}
\usepackage[unicode,hidelinks,bookmarks=false]{hyperref}
\newcommand{\ie}{i.e.\ }
\newcommand{\cf}{cf.\ }
\newcommand{\resp}{respectively\ }
\newcommand{\Subsection}{\subsection}
\providecommand{\nobreakdash}{\nobreak\hbox{-}}
\setlength{\emergencystretch}{2em}
\multlinegap0pt
\usepackage{amssymb}
\usepackage{enumerate}
\newtheorem{lem}{Lemma}
\title{The Grünbaum--Rigby configuration as a special Kárteszi configuration}
\author{G\'abor G\'evay, Gy\"orgy Kiss, Toma\v{z} Pisanski}
\date{Source: arXiv:2512.18872v1 (CC BY 4.0)}
\begin{document}
\maketitle

\noindent\emph{Source and licence.} The Grünbaum--Rigby configuration as a special Kárteszi configuration. Version arXiv:2512.18872v1 (CC BY 4.0), distributed by arXiv under the Creative Commons Attribution 4.0 International licence, http://creativecommons.org/licenses/by/4.0/, as recorded in the arXiv licence metadata for this exact version and shown on the abstract page https://arxiv.org/abs/2512.18872v1. Full licence notice: \texttt{licenses/arxiv-2512.18872-CC-BY-4.0.txt}.\par\smallskip

\noindent\textit{Editorial scope.} Counted unit: the article's lem 3-bol4 (source0.tex lines 448--460). The article's complete original proof of it is delivered with it (source0.tex lines 462--478, one \texttt{proof} environment of 17 lines). No mathematics is written, repaired or re-derived: the statement and the proof are the article's own lines, in the article's own notation, and no retained line is substituted or re-flowed.  No further source-local context is needed: every cross-reference of every retained line resolves inside the two delivered spans.  Nothing was omitted from any retained span, so the package carries no omission record.  No retained line contains a citation, so the wrapper carries no bibliography and no bibliography entry.  No retained line contains a figure environment, an image-inclusion command or drawing code, so no image is delivered and the package declares no asset dependency.  Editorial formatting only, in addition: an AMS \texttt{amsart} preamble that restates the article's own declarations and macros used by the retained lines, namely the environments \texttt{lem} on separate counters, together with the standard packages those lines need.  The retained environments print in this document as Lemma 1 in order of appearance, whereas the article prints the same items with the numbering of its own full document.  The complete native source of the article as distributed in the arXiv e-print for this exact version is bundled unchanged as \texttt{sources/191427/source0.tex}.
\begin{lem}
\label{3-bol4}
Suppose that an $r$-th diagonal $d$ of $\mathcal{A}_1 $ is incident with a vertex $B$ of $\mathcal{A}_k$. 
If $B$ is different from the midpoint of $d$, then $d$ is incident with another vertex of $\mathcal{A}_k $, 
hence $d$ has a common line with a diagonal of $\mathcal{A}_k $.

Let $B_i=A_iA_{i+k}\cap A_{i+1}A_{i+1+k}$ denote the vertices of
$\mathcal{A}_k $ (the subscripts are calculated modulo $n$). 
If the diagonal $A_0A_{r}$ is incident with $B_i$, $r<\frac{n}{2},$ $\frac{r}{2}<i$ and 
$2i-r+k+1<\frac{n}{2}$, then $B_{r-i-k-1}$ is also incident with $A_0A_{r}$. 
The $r$-th diagonals of $\mathcal{A}_1 $
have a common line with the $(2i-r+k+1)$-th diagonals of $\mathcal{A}_k $. 
\end{lem}

\begin{proof}
Let $t$ be the perpendicular bisector of $d$. Then $t$ is an axis of symmetry of $\mathcal{A}_1 $, 
hence it is also an axis of symmetry of $\mathcal{A}_k $. Thus, if $\varphi$ denotes the reflection 
in $t$, then $\varphi (B)\neq B$ is another vertex of $\mathcal{A}_k $ and $d$ has a common line 
with the diagonal  $B\varphi (B)$ of $\mathcal{A}_k $.

If $d$ is the diagonal $A_0A_{r}$, then we have $\varphi (A_j)=A_{r-j}$ 
for $j=0,1,\dots ,n-1.$ So 
$$\varphi (B_i)=\varphi (A_{i}A_{i+k}\cap A_{i+1}A_{i+1+k}) =A_{r-i}A_{r-i-k}
\cap A_{r-i-1}A_{r-i-k-1}.$$ 
Since $i+1+k< n-(r-i-k-1)$ we get $A_{r-i-1}A_{r-i-k-1}=B_{r-k-i-1}$. 
Thus, $d$ is incident with $B_{r-k-i-1}$, so it has a common  line with the diagonal 
$B_iB_{r-k-i-1}$ of $\mathcal{A}_k $. This diagonal is a $(2i-r+k+1)$-th diagonal, 
because $2i-r+k+1<n/2.$ By rotational symmetry, this implies that any
$r$-th diagonal of $\mathcal{A}_1 $
has a common line with a suitable $(2i-r+k+1)$-th diagonal of $\mathcal{A}_k $. 
\end{proof}

\end{document}
